\documentclass[a4paper,10pt,twocolumn,oneside]{article}

\usepackage{amsmath, amssymb, amsthm}    % 美國數學學會
\usepackage[CheckSingle, CJKmath]{xeCJK} % 中文
\usepackage{fontspec}                    % 字型配置
\usepackage{geometry}                    % 版面配置
\usepackage{fancyhdr}                    % 頁首頁尾
\usepackage{color}                       % 顏色
\usepackage[x11names]{xcolor}            % 更多顏色
\usepackage{listings}                    % 顯示 code 用的
\usepackage{tikz}                        % 畫圖
\usepackage{enumerate}
\usepackage{graphicx}
\usepackage{multicol}
\usepackage{titlesec}
\usepackage{xpatch}
% \usepackage{courier}
% \usepackage[Glenn]{fncychap}           % 排版，頁面模板
% \usepackage{CJKulem}
% \usepackage[T1]{fontenc}

%%%%%%%%%%%%%%%%%%%%%%%%%%%%%

\setmainfont{NotoSans}[  % 主要字體
    Path           = .fonts/Noto_Sans/,
    Extension      = .ttf,
    UprightFont    = *-Regular,
    BoldFont       = *-Bold,
    ItalicFont     = *-Italic,
    BoldItalicFont = *-BoldItalic,
]

\setmonofont{Consolas}[  % 等寬字體
    Path           = .fonts/Consolas/,
    Extension      = .ttf,
    UprightFont    = *-Regular,
    BoldFont       = *-Bold,
    ItalicFont     = *-Italic,
    % BoldItalicFont = *-BoldItalic,
]

\setCJKmainfont{NotoSansTC}[  % 中文主要字體
    Path           = .fonts/Noto_Sans_TC/,
    Extension      = .ttf,
    UprightFont    = *-Regular,
    BoldFont       = *-Bold,
]

% \setCJKmonofont{LXGWWenKaiMonoTC}[  % 中文等寬字體
% Path           = .fonts/LXGW_WenKai_Mono_TC/,
% Extension      = .ttf,
% UprightFont    = *-Regular,
% BoldFont       = *-Bold,
% ]

%%%%%%%%%%%%%%%%%%%%%%%%%%%%%

% 彩色模式
\newcommand{\keywordcolor}{\color{Blue1}}
\newcommand{\identifiercolor}{\color{black}}
\newcommand{\commentcolor}{\color{Red4}}
\newcommand{\stringcolor}{\color{Green4}}

% 黑白模式
% \renewcommand{\keywordcolor}{\color{black}}
% \renewcommand{\identifiercolor}{\color{black}}
% \renewcommand{\commentcolor}{\color{black!70}}
% \renewcommand{\stringcolor}{\color{black!70}}

\makeatletter
\lst@CCPutMacro\lst@ProcessOther {"2D}{\lst@ttfamily{-{}}{-{}}}
\@empty\z@\@empty
\makeatother

\lstset{
    language=C++,                           % Code 的語言
    numbers=none,                           % 行號的位置
    numberstyle=\footnotesize,              % 行號的字型和大小
    stepnumber=1,                           % 顯示行號的間隔
    numbersep=5pt,                          % 行號跟 code 的距離
    showspaces=false,                       % 顯示空白 (使用特別的底線記號)
    showtabs=false,                         % 顯示 Tab
    showstringspaces=false,                 % 顯示字串的空白
    tabsize=2,                              % Tab 的寬度
    frame=leftline,                         % adds a frame around the code
    captionpos=b,                           % sets the caption-position to bottom
    breaklines=true,                        % sets automatic line breaking
    breakatwhitespace=false,                % sets if automatic breaks should only happen at whitespace
    escapeinside={\%*}{*)},                 % if you want to add a comment within your code
    morekeywords={*},                       % 自訂的 Keywords
    basicstyle=\footnotesize\ttfamily,      % 基本的字型和大小
    backgroundcolor=\color{white},          % 背景顏色
    keywordstyle=\bfseries\keywordcolor,    % Keywords 的字型
    identifierstyle=\identifiercolor,       % 標示符的字型
    commentstyle=\itshape\commentcolor,     % 註解的字型
    stringstyle=\itshape\stringcolor,       % 字串的字型
}

%%%%%%%%%%%%%%%%%%%%%%%%%%%%%

\geometry{a4paper}                           % 頁面 A4 大小
\geometry{includehead, nofoot, nomarginpar}  % 內文區塊包含頁首，沒有頁尾、旁注
\geometry{headsep=2mm}                       % 頁首和內文的距離
\geometry{vmargin=5mm, hmargin=5mm}          % 垂直、水平邊界

\setlength{\columnsep}{3mm}                  % 兩欄模式的間距
\setlength{\columnseprule}{0pt}              % 兩欄模式間格線粗細

\titlespacing{\section}{0cm}{0cm}{0cm}
\titlespacing{\subsection}{0cm}{0cm}{0cm}

\setcounter{secnumdepth}{3}                  % 目錄顯示第三層

\XeTeXlinebreaklocale "zh"                   % 中文自動換行
\XeTeXlinebreakskip = 0pt plus 1pt           % 設定段落之間的距離

%%%%%%%%%%%%%%%%%%%%%%%%%%%%%

\begin{document}
\pagestyle{fancy}

%%%%%

% 頁首、頁尾
\renewcommand{\headrulewidth}{0.4pt}
\fancyhead[L]{NTOU: Haruhikagerogerokonokositantan (ﾊﾙﾋｶｹﾞﾛｹﾞﾛｺﾉｺｼﾀﾝﾀﾝ)}
\fancyhead[R]{\thepage}
\fancyfoot{}

%%%%%

% 目錄
\renewcommand{\contentsname}{Contents}
\scriptsize
\begin{multicols}{2}
    \tableofcontents
\end{multicols}

%%%%%

% 欄寬測試
% \lstinputlisting{misc/width_test.txt}

%%%%%%%%%%%%%%%%%%%%%%%%%%%%%

\section{Basis}


\subsection{Default Code}
\lstinputlisting{basis/default_code.cpp}

\subsection{Some Useful Primes}
\lstinputlisting{basis/primes.cpp}

\subsection{隨機 \& MT19937}
\lstinputlisting{basis/mt.cpp}

%%%%%

\section{Data Structures}

\subsection{Disjoint Set (Union-Find)}
\lstinputlisting{data_structure/dsu.cpp}

\subsection{Binary Indexed Tree}
\lstinputlisting{data_structure/binary_indexed_tree.cpp}

\subsection{Segment Tree}
\lstinputlisting{data_structure/segment_tree.cpp}

%%%%%

\section{Mathematics}

\subsection{Combination with Inverse}
\lstinputlisting{math/combinations_with_inverse.cpp}

\subsection{EXGCD}
\lstinputlisting{math/exgcd.cpp}

\subsection{FFT}
\lstinputlisting{math/fast_fourier_transform.cpp}

\subsection{Exponentiation by Squaring}
\lstinputlisting{math/fast_pow.cpp}

\subsection{GCD and LCM}
\lstinputlisting{math/gcd_lcm.cpp}

\subsection{Miller-Rabin Algorithm}
\lstinputlisting{math/miller_rabin.cpp}

\subsection{Pollard's Rho Algorithm}
\lstinputlisting{math/pollard_rho.cpp}

%%%%%

\section{Graph}

\subsection{Topological Sort}
\lstinputlisting{graph_theory/topological_sort.cpp}

\subsection{Tarjan's Algorithm for BCC}
\lstinputlisting{graph_theory/bcc_tarjan.cpp}

\subsection{Bellmen-Ford Algorithm}
\lstinputlisting{graph_theory/bellmen_ford.cpp}

\subsection{Dijkstra's Algorithm}
\lstinputlisting{graph_theory/dijkstra.cpp}

\subsection{Dinic's Algorithm for Maximum Flow}
\lstinputlisting{graph_theory/dinic.cpp}

\subsection{Dominator Tree}
\lstinputlisting{graph_theory/dominator_tree.cpp}

\subsection{Floyd-Warshall Algorithm}
\lstinputlisting{graph_theory/floyd_warshall.cpp}

\subsection{Maximum Clique}
\lstinputlisting{graph_theory/maximum_clique.cpp}

\subsection{Kosaraju's Algorithm for SCC}
\lstinputlisting{graph_theory/scc_kosaraju.cpp}

%%%%%

\section{Tree Algorithms}

\subsection{Lowest Common Ancester}
\lstinputlisting{tree_algorithm/lowest_common_ancester.cpp}

\subsection{Eular Tour for LCA}
\lstinputlisting{tree_algorithm/eular_tour_for_lca.cpp}

\subsection{DSU on Tree}
\lstinputlisting{tree_algorithm/dsu_on_tree.cpp}

%%%%%

\section{String Algorithms}

\subsection{Suffix Array (SAIS)}
\lstinputlisting{string/SAIS.cpp}

\subsection{Z-algorithm}
\lstinputlisting{string/z_algorithm.cpp}

\subsection{Knuth-Morris-Pratt Algorithm}
\lstinputlisting{string/KMP.cpp}

\subsection{Manacher's Algorithm}
\lstinputlisting{string/manacher.cpp}

\subsection{Minimum Rotation}
\lstinputlisting{string/min_rotation.cpp}

\subsection{Rolling Hash}
\lstinputlisting{string/rolling_hash.cpp}

%%%%%

\section{Geometry}

\subsection{Basic Geometry Structure}
\lstinputlisting{geometry/geometry_struct.cpp}

\subsection{Circle Cover Area}
\lstinputlisting{geometry/circle_cover.cpp}

\subsection{Circles' Tangent Lines}
\lstinputlisting{geometry/circle_tangent_line.cpp}

\subsection{Circles' Intersections}
\lstinputlisting{geometry/circle_intersection.cpp}

\subsection{Convex Hull}
\lstinputlisting{geometry/convex_hull.cpp}

\subsection{Convex Hull Tricks}
\lstinputlisting{geometry/convex_hull_tricks.cpp}

\subsection{Pick's Theory}
\lstinputlisting{geometry/picks_theory.cpp}

%%%%%

\section{Sorting}

\subsection{Merge Sort}
\lstinputlisting{sort/merge_sort.cpp}

\subsection{Quick Sort}
\lstinputlisting{sort/quick_sort.cpp}

%%%%%

\section{Miscellaneous}

\subsection{Discretization}
\lstinputlisting{misc/discretization.cpp}

\subsection{Blessing}
\lstinputlisting{misc/blessing1.cpp}
\lstinputlisting{misc/blessing2.cpp}
\lstinputlisting{misc/blessing3.cpp}

%%%%%

\clearpage

\subsection{Graph Paper}

% 產生方格紙 \graph{X 長度}{Y 長度}
\newcommand{\graph}[2]{
    \noindent
    \begin{tikzpicture}[every node/.style={minimum size=1cm-\pgflinewidth, outer sep=10pt}, scale=2]
        \draw[step=0.2cm, color=gray] (0, 0) grid (#1, #2);
        \draw[step=1cm, color=black, line width=0.4mm] (0, 0) grid (#1, #2);
    \end{tikzpicture}
}

\graph{10}{13.6}
\clearpage
\graph{10}{14}
\clearpage
\graph{10}{14}
\clearpage
% \graph{10}{14}
% \clearpage

\end{document}
